% ============================================================
%  Seção 2.4 — Análise de Sentimento
%  Capítulo 2 — Fundamentação Teórica
% ============================================================

\section{Análise de Sentimento}
\label{sec:analise_sentimento}

A Análise de Sentimento (AS), também denominada mineração de opinião (\textit{opinion mining}), é a área do PLN dedicada à identificação e extração computacional de opiniões, sentimentos, emoções e atitudes expressas em textos não estruturados \cite{liu2020survey}. O objetivo central da AS é determinar a \textbf{polaridade} de um dado conteúdo --- se ele expressa uma orientação positiva, negativa ou neutra --- podendo operar em diferentes níveis de granularidade: no nível do documento, da sentença ou de aspectos específicos do texto \cite{jurafsky2026,liu2020survey}. No domínio financeiro, a AS adquire relevância particular ao permitir mensurar a percepção coletiva dos agentes econômicos --- expressa em notícias, relatórios e publicações em redes sociais --- fornecendo um indicador qualitativo do humor do mercado que complementa as métricas quantitativas tradicionais \cite{kearney2014,tetlock2007}.

% ============================================================
%  Subseção 2.4.1 — Tarefas e Níveis de Análise
% ============================================================

\subsection{Tarefas e Níveis de Análise}
\label{subsec:as_tarefas}

A AS pode ser organizada em torno de três tarefas principais, que diferem em complexidade e granularidade \cite{liu2020survey,jurafsky2026}. A \textbf{classificação de	polaridade} é a tarefa mais diretamente relevante para este trabalho: dado um texto, o sistema atribui-lhe uma categoria de sentimento --- positivo, negativo ou neutro. Em contextos financeiros, essa classificação frequentemente se traduz em sinais de otimismo ou pessimismo do investidor em relação a um ativo ou ao mercado como um todo \cite{bollen2011,galdi2018}. A \textbf{análise de sentimento baseada em aspectos} (\textit{aspect-based sentiment analysis}) vai além da polaridade global e identifica o sentimento associado a entidades ou atributos específicos mencionados no texto --- por exemplo, distinguindo a avaliação sobre a gestão de uma empresa de sua avaliação sobre seus resultados financeiros \cite{liu2020survey}. A \textbf{detecção de emoções} vai além da polaridade binária e classifica o conteúdo emocional em categorias como alegria, medo, raiva e tristeza, conforme modelos psicológicos de afeto \cite{bollen2011}; essa dimensão foi explorada de forma pioneira por \citeauthor{bollen2011} (\citeyear{bollen2011}) ao decompor o humor coletivo do Twitter em seis dimensões do GPOMS (\textit{Google-Profile of Mood States}) e correlacioná-las ao DJIA.

No que diz respeito aos níveis de análise, a literatura distingue entre a análise em nível de \textbf{documento} --- em que uma publicação ou notícia inteira é classificada com uma única polaridade --- e a análise em nível de \textbf{sentença} --- em que cada sentença é classificada individualmente, permitindo capturar a coexistência de sentimentos contraditórios num mesmo texto \cite{kearney2014}. Para publicações curtas como \textit{tweets}, os dois níveis tendem a coincidir, dada a natureza unitária e concisa de cada publicação \cite{bollen2011,ranco2015}. Este trabalho opera predominantemente no nível de documento/publicação, classificando cada \textit{tweet} individualmente como positivo, negativo ou neutro.

% ============================================================
%  Subseção 2.4.2 — Abordagens Metodológicas
% ============================================================

\subsection{Abordagens Metodológicas}
\label{subsec:as_abordagens}

As abordagens para realização de AS podem ser organizadas em três grandes famílias, cada uma com características, vantagens e limitações distintas \cite{kearney2014,liu2020survey,lucas2025}.

A \textbf{abordagem baseada em léxicos} classifica o sentimento de um texto computando a soma (ou média ponderada) das pontuações de sentimento das palavras individuais que o compõem, conforme um dicionário pré-compilado que associa cada palavra a um valor de polaridade \cite{loughran2011,sousa2025}. As principais vantagens dessa abordagem são a interpretabilidade, a ausência de necessidade de dados rotulados para treinamento e a facilidade de incorporação de conhecimento especializado de domínio \cite{lucas2025}. Suas limitações são igualmente conhecidas: léxicos de propósito geral falham em capturar o significado específico do vocabulário financeiro, conforme demonstrado por \citeauthor{loughran2011} (\citeyear{loughran2011}); a abordagem é sensível a fenômenos de negação (e.g., ``\textit{não subiu}''), ironia e sarcasmo; e não captura dependências contextuais entre palavras distantes no texto \cite{liu2020survey}. No contexto brasileiro, \citeauthor{lucas2025} (\citeyear{lucas2025}) demonstraram que os léxicos VADER e HIV conseguem explicar, respectivamente, 28,8\% e 18,4\% da variância do Ibovespa quando aplicados a notícias financeiras em português --- resultado significativo, mas limitado pela inadequação dos léxicos ao vocabulário do mercado brasileiro.

A \textbf{abordagem baseada em aprendizado de máquina} treina modelos estatísticos ou neurais sobre conjuntos de textos rotulados para aprender, a partir dos dados, os padrões linguísticos associados a cada classe de sentimento \cite{liu2020survey}. Métodos tradicionais como Naïve Bayes, Máquinas de Vetores de Suporte (SVM) e árvores de decisão operam sobre representações esparsas do texto (\textit{bag-of-words}, TF-IDF) e alcançam bom desempenho em domínios com vocabulário estável e dados rotulados suficientes \cite{sousa2025}. A principal limitação dessa família é a dependência de grandes corpora anotados, que são escassos no domínio financeiro em português \cite{santos2023finbert,sousa2025}. Além disso, modelos treinados em textos de um domínio ou gênero específico generalizam com dificuldade para outros contextos.

A \textbf{abordagem baseada em modelos de linguagem profundos}, em particular os modelos \textit{transformer} discutidos na Seção~\ref{subsec:transformers}, representa o estado da arte atual em AS \cite{liu2020survey,santos2023finbert}. Esses modelos aprendem representações contextualizadas das palavras e sentenças durante o pré-treinamento em grandes corpora não rotulados e são depois ajustados (\textit{fine-tuned}) para a tarefa de classificação de sentimento com um número relativamente pequeno de exemplos rotulados. A contextualização é a chave para superar as limitações das abordagens anteriores: o significado de cada palavra é função de todo o contexto em que ocorre, permitindo que o modelo capture negações, ambiguidades e nuances semânticas que os léxicos e os modelos \textit{bag-of-words} não conseguem representar \cite{jurafsky2026,devlin2019}. A \textbf{abordagem híbrida}, por fim, combina elementos léxicos com modelos de aprendizado, aproveitando a especificidade dos dicionários especializados para enriquecer as representações aprendidas pelo modelo neural \cite{kearney2014,sousa2025}.

% ============================================================
%  Subseção 2.4.3 — Análise de Sentimento no Domínio Financeiro
%                   em Português
% ============================================================

\subsection{Análise de Sentimento no Domínio Financeiro em Português}
\label{subsec:as_financeiro_ptbr}

A aplicação da AS ao domínio financeiro em língua portuguesa apresenta desafios e oportunidades específicos que justificam o esforço de especialização de recursos e modelos. Como discutido na Seção~\ref{subsec:desafios_textuais}, a transferência direta de ferramentas desenvolvidas para o inglês financeiro produz resultados sistematicamente inferiores ao potencial da área \cite{loughran2011,santos2023finbert}. O vocabulário do mercado financeiro brasileiro é rico em termos técnicos, siglas de índices e ativos (Ibovespa, IPCA, Selic, PETR4), referências institucionais (B3, CVM, Banco Central) e expressões idiomáticas que não possuem correspondentes diretos em léxicos ou corpora em inglês \cite{sousa2025,falcao2020}.

Os primeiros trabalhos aplicados ao mercado brasileiro utilizaram principalmente abordagens léxicas e modelos de aprendizado de máquina raso. \citeauthor{falcao2020} (\citeyear{falcao2020}) analisou notícias do Valor Econômico e da Folha de São Paulo entre 2013 e 2019, correlacionando o sentimento extraído com o comportamento do Ibovespa e de cinco ações de diferentes setores, e constatou que os sentimentos tendem a ser neutros na maioria dos dias, seguindo tendências pessimistas em períodos de maior incerteza econômica. \citeauthor{galdi2018} (\citeyear{galdi2018}) investigaram a relação entre pessimismo e incerteza das notícias e o comportamento dos investidores brasileiros, evidenciando que o conteúdo informacional da mídia influencia a visão dos investidores especialmente em momentos de maior turbulência do mercado.

A consolidação do paradigma \textit{transformer} no PLN abriu novas perspectivas para o domínio financeiro em português. O modelo FinBERT-PT-BR \cite{santos2023finbert}, desenvolvido por \citeauthor{santos2023finbert} (\citeyear{santos2023finbert}) a partir do ajuste fino do BERTimbau \cite{souza2020} sobre 1,4 milhão de sentenças de notícias financeiras em português, estabeleceu um novo estado da arte para a tarefa de classificação de sentimento no domínio. O modelo é publicamente disponível e foi avaliado em um
conjunto de 500 sentenças rotuladas de notícias financeiras, superando modelos multilíngues e léxicos especializados nas métricas de acurácia, F1-Score, precisão e recall \cite{santos2023finbert}. Iniciativas paralelas de construção de léxicos especializados para o português financeiro, como a proposta de expansão automática por Word2Vec e PMI de \citeauthor{sousa2025} (\citeyear{sousa2025}), alcançaram F1-Score de 80\% em \textit{tweets} quando combinadas com classificadores SVM, evidenciando o potencial das abordagens híbridas para complementar os modelos neurais.

No presente trabalho, adota-se o modelo FinBERT-PT-BR como ferramenta central para a extração de polaridade das publicações do X/Twitter relacionadas ao mercado financeiro brasileiro. Essa escolha é motivada por três fatores: (i) a especialização do modelo no vocabulário financeiro em português do Brasil; (ii) seu desempenho superior ao estado da arte anterior nas tarefas de classificação de sentimento financeiro; e (iii) sua disponibilidade pública, que garante a reprodutibilidade da metodologia proposta. Complementarmente, o trabalho avalia o BERTimbau \cite{souza2020} ajustado por \textit{fine-tuning} diretamente sobre o corpus de \textit{tweets} financeiros coletados, como termo de comparação que situa o desempenho do FinBERT-PT-BR em relação a um modelo genérico especializado \textit{a posteriori} para o gênero e o subdomínio específicos do corpus --- abordagem detalhada na \autoref{subsec:bertimbau}. A aplicação de ambos os modelos ao corpus coletado, bem como a análise dos índices de sentimento resultantes, é detalhada no Capítulo~\ref{Cap:Metodologia}.